% titel: Blatt 2, A8: Das Haber-Bosch-Verfahren
% Teil II: die erarbeiteten Antworten.

\subsection*{Fragen und Antworten}
\begin{itemize}
  \item[a)] \textbf{Reaktionsgleichung des Haber-Bosch-Verfahrens:}
        \begin{equation}
          \mathrm{N_2} + 3\,\mathrm{H_2} \;\rightleftharpoons\;
          2\,\mathrm{NH_3}
        \end{equation}
        Die Reaktion ist eine Gleichgewichtsreaktion (reversible Form), die
        technisch über einen Katalysator in Richtung Produkt gefahren wird.
  \item[b)] \textbf{Begrenzender Reaktand:} In der Regel ist
        $\mathrm{H_2}$ der begrenzende Reaktand, weil Stickstoff aus der Luft
        (fast unbegrenzt) verfügbar ist, waehrend Wasserstoff aufwändig
        (z.~B. aus Methan) hergestellt werden muss und in begrenzter Menge
        dosiert wird.
  \item[c)] \textbf{Anwendung:} Das Verfahren dient der Herstellung von
        Ammoniak $\mathrm{(NH_3)}$, überwiegend für Düngemittel.
  \item[d)] \textbf{Substrat vs.\ Edukt:} \emph{Edukt} ist jeder
        Ausgangsstoff einer chemischen Reaktion. \emph{Substrat} nennt man
        denjenigen Stoff, der mithilfe eines Katalysators umgewandelt wird;
        der Begriff wird also vor allem bei katalysierten Reaktionen
        verwendet.
\end{itemize}

\begin{ergebnis}
\begin{align*}
  \mathrm{N_2} + 3\,\mathrm{H_2} &\to 2\,\mathrm{NH_3}
    \quad(\text{Haber-Bosch}),\\
  \text{begrenzender Reaktand} &= \mathrm{H_2}\quad(\text{regelmässig}),\\
  \text{Zweck} &= \text{Ammoniaksynthese (Düngemittel)},\\
  \text{Edukt} &= \text{allgemeiner Ausgangsstoff},\quad
  \text{Substrat} = \text{katalysiert umgesetzter Stoff}.
\end{align*}
\end{ergebnis}

\begin{kontrolle}
Die Gleichung ist stochiometrisch konsistent (Atomerhalt: $2$ N und $6$ H je
Seite). Die Aussage zum begrenzenden Reaktanden folgt aus den tatsächlich
eingesetzten Stoffmengen in der Technik (H$_2$ knapp, N$_2$ reichlich aus der
Luft).
\end{kontrolle}

\begin{merksatz}
Industriesynthesen laufen über ein ausbalanciertes Gleichgewicht; der
begrenzende Reaktand (hier meist H$_2$) bestimmt, wie viel Produkt maximal
entsteht, und ``Substrat'' ist der katalytisch umgesetzte Stoff im Gegensatz
zum Oberbegriff ``Edukt''.
\end{merksatz}
